The correction of handwritten exams is a time-intensive process that
lacks platforms tailored to the Spanish academic context. Existing
solutions, mostly foreign, present language barriers and incomplete
integration with national academic management systems.

This Bachelor's Thesis presents the design and implementation of the
\dfn{backend} of SCDEE (Sistema de Corrección Digital de Exámenes
Escritos — Digital Handwritten Exam Correction System), a distributed
\dfn{multi-tenant} service that automates the complete correction cycle,
from the ingestion of scanned pages to the final archival of the
corrected exam. The system integrates a recognition \textit{pipeline}
based on \acs{qr} codes and error-tolerant \acs{ocr}, built on a modular
architecture of thirteen Django applications with a service layer
extensible through interchangeable \dfnpl{plugin} (authentication,
\acs{ocr} engines, zone recognisers).

The \dfn{backend} implements \dfn{multi-tenant} access control with
contextual permissions per subject, an intelligent exam-to-student
assignment mechanism, and both manual and rubric-based grading. Sensitive
data is protected through field-level encryption, robust password
hashing, and dynamic watermarks on the \acs{pdf} files served to
students, all backed by an immutable audit trail of every operation. The
system also offers per-organisation configuration of recognition
parameters, limits and permissions, as well as automated maintenance
tasks with email notifications.

System validation combines close to 90\,\% coverage in unit and
integration tests, load tests with Locust that verify
\cref{rnf:concurrency} under normal operation, and an empirical
evaluation of \acs{ocr} quality that justifies the choice of the EasyOCR
engine. The architecture is prepared to evolve toward production-grade
deployments, integration with \acs{lms} systems and institutional
identity federation.

\keywords{digital exam correction, multi-tenancy, OCR, Django,
distributed systems, backend architecture}
